\section{Design Definition and Analytical Baseline}
\label{sec:design-definition}

\subsection{Geometry and Operating Conditions}

The modeled vacuum region has the dimensions

\[
a = 20~\mathrm{mm},\qquad
b = 20~\mathrm{mm},\qquad
l = 90~\mathrm{mm}.
\]

The transverse dimensions are equal, so the structure is a square waveguide within the broader rectangular-waveguide family. The frequency-domain simulation covers

\[
8~\mathrm{GHz}\le f \le 15~\mathrm{GHz},
\]

and the principal field-evaluation frequency is

\[
f_0 = 10~\mathrm{GHz}.
\]

\begin{table}[H]
\centering
\caption{Public model definition.}
\label{tab:model-definition}
\begin{tabular}{@{}lll@{}}
\toprule
\textbf{Quantity} & \textbf{Value} & \textbf{Model interpretation} \\
\midrule
Width \(a\) & \(20~\mathrm{mm}\) & First transverse dimension \\
Height \(b\) & \(20~\mathrm{mm}\) & Second transverse dimension \\
Length \(l\) & \(90~\mathrm{mm}\) & Uniform propagation section \\
Interior & Vacuum & Lossless filling medium \\
Walls & PEC & Ideal conductor, zero conductor loss \\
Frequency sweep & \(8\)--\(15~\mathrm{GHz}\) & Frequency-domain solution \\
Field reference & \(10~\mathrm{GHz}\) & Modal field comparison \\
Port modes & 3 per port & Lowest three solved eigenmodes \\
\bottomrule
\end{tabular}
\end{table}

\subsection{Why the Square Cross Section Matters}

For a conventional rectangular horn feed, one normally selects \(a>b\) so that the \(\mathrm{TE}_{10}\) mode has a lower cutoff frequency than \(\mathrm{TE}_{01}\) \cite{pozar2011microwave,stutzman2013antenna}. Here, \(a=b\), which forces

\[
f_{c,10}=f_{c,01}.
\]

The model therefore supports two orthogonally polarized lowest-order modes at the same cutoff frequency. A numerical port solver may choose any orthogonal basis spanning that degenerate subspace. The selected basis can differ between Port 1 and Port 2, so mode index alone is not a reliable physical polarization label \cite{pozar2011microwave}.

\subsection{Expected Behavior at 10 GHz}

Using the exact vacuum speed of light, the first two cutoff frequencies are expected near \(7.4948~\mathrm{GHz}\). The \((1,1)\)-family cutoff is expected near \(10.5993~\mathrm{GHz}\). Therefore:

\begin{itemize}
  \item the first two solved modes should propagate at \(10~\mathrm{GHz}\);
  \item the third solved mode should be slightly below cutoff at \(10~\mathrm{GHz}\);
  \item the first two port modes should exhibit nearly identical propagation constants and wave impedances;
  \item the output modal amplitudes may be rotated between the two degenerate basis functions.
\end{itemize}
